% Auto-generated — REAL cluster timings (mode=cluster).
% Highlight: N=5 / C1  runs=1

\subsection{Orthogonal Testing Efficiency}

\begin{figure*}[t]
\centering
% \includegraphics[width=0.85\linewidth]{figures/evaluation/orthogonal-reduction.png}
\caption{Orthogonal validation under different scales (real cluster).
(a) Test-case growth vs configuration items.
(b) Complete validation wall-clock time vs configuration items.}
\label{fig:orthogonal-reduction}
\end{figure*}

\textbf{Goal.}
Evaluate orthogonal reduction of dynamic validation cost on a live Istio cluster, and whether fault detection coverage is preserved after merging.

\textbf{Setup.}
Service scales $N\in\{5\}$, configuration levels $\{C1\}$.
Each validation case performs live \texttt{kubectl apply}, ingress traffic, assertion, and delete.
Orthogonal bundles share one apply/clear and issue per-effect traffic.
Each cell is repeated 1 times; times are wall-clock mean$\pm$std.

\textbf{Results.}
Under a 5-service deployment at C1 (70 configuration items), exhaustive validation uses 25 cases vs 7 orthogonal cases (reduction 72.00\%).
Complete validation wall-clock time drops from $19.1\pm0.0$s to $5.5\pm0.0$s (time reduction 71.34\%).

\textbf{Coverage ablation.}
Table~\ref{tab:ortho-ablation} reports live fault/effect detection before vs after merging.

\begin{table*}[t]
\centering
\caption{Real-cluster detection and validation time (C1).}
\label{tab:ortho-ablation}
\scriptsize
\begin{tabular}{cccccccccc}
\toprule
$N$ & $|F|$ & DR$_b$ & DR$_a$ & $\Delta$ & Cases$_e$ & Cases$_o$ & Red. & $T_e$ (s) & $T_o$ (s) \\
\midrule
5 & 25 & 100.00\% & 100.00\% & 0.00 & 25 & 7 & 72.00\% & 19.1$\pm$0.0 & 5.5$\pm$0.0 \\
\bottomrule
\end{tabular}
\end{table*}
